% lecture10.tex: Lecture 10: Rocky Planets II, Mercury & Mars
% Planetary Systems, Kapteyn Institute

\documentclass[aspectratio=169, 11pt]{beamer}

% ── Theme and macros ────────────────────────────────────────
\usepackage{../common/beamerthemeIPS}
% macros.tex: Shared math macros for IPS lecture slides
% Mirrors definitions in book/_config.yml for consistency
% Uses \providecommand to avoid conflicts with unicode-math

% ── Derivatives ─────────────────────────────────────────────
\providecommand{\dv}[2]{\frac{\mathrm{d} #1}{\mathrm{d} #2}}
\providecommand{\dd}{}\renewcommand{\dd}{\mathrm{d}}
\providecommand{\pdv}[2]{\frac{\partial #1}{\partial #2}}

% ── Solar quantities ────────────────────────────────────────
\newcommand{\Msun}{M_{\odot}}
\newcommand{\Rsun}{R_{\odot}}
\newcommand{\Lsun}{L_{\odot}}

% ── Planetary quantities ────────────────────────────────────
\newcommand{\Mearth}{M_{\oplus}}
\newcommand{\Rearth}{R_{\oplus}}
\newcommand{\Mjup}{M_{\mathrm{J}}}
\newcommand{\Rjup}{R_{\mathrm{J}}}

% ── Constants ───────────────────────────────────────────────
\newcommand{\kB}{k_{\mathrm{B}}}

% ── Media links ─────────────────────────────────────────────
% Clickable button that opens the animated or video version of a
% figure, e.g. \mediabutton{https://...gif}{Play animation}
\providecommand{\mediabutton}[2]{\href{#1}{\beamergotobutton{#2}}}

\graphicspath{{figures/}{../common/}{../../book/10_mercury_mars/figures/}}

% ── Metadata ────────────────────────────────────────────────
\title[Lecture 10: Mercury \& Mars]{Lecture 10: Rocky Planets II\\Mercury \& Mars}
\subtitle{Planetary Systems}
\author[L10]{Tim Lichtenberg}
\institute{Kapteyn Astronomical Institute\\University of Groningen}
\date{2026}
\titlebgcredit{Mercury topography from MESSENGER MLA. Credit: NASA/JHUAPL/CIW}

% ════════════════════════════════════════════════════════════
\begin{document}

% ── Title slide ─────────────────────────────────────────────
\begin{frame}[plain,noframenumbering]
  \titlepage
\end{frame}

% ── Learning objectives ─────────────────────────────────────
\begin{frame}{Learning Objectives}
  By the end of this lecture, you will be able to:
  \vskip8pt
  \begin{itemize}
    \item Explain Mercury's large iron core, its 3:2 spin-orbit resonance and its weak offset dynamo
    \item Describe Mars' interior from the InSight seismic results and its wet-to-dry surface record
    \item Derive the Jeans escape flux, apply it to H, $\mathrm{H_2}$ and $\mathrm{CO_2}$ on Mars, and distinguish thermal from non-thermal escape
  \end{itemize}
\end{frame}

% ── Recap ───────────────────────────────────────────────────
\begin{frame}{Recap: Lecture 9}
  \begin{columns}[c]
    \column{\recapfigcol}
    \ipsrecapfig{terrestrial_planets}
    \source{NASA/JPL}

    \column{\recaptextcol}
    \begin{itemize}
      \item Earth and Venus: similar bulk parameters, radically divergent evolution
      \item Simpson-Nakajima runaway limit: maximum OLR of a saturated water atmosphere $\sim 280$ W\,m$^{-2}$
      \item Plate tectonics, dynamo, and biosphere form an Earth-only coincidence
    \end{itemize}
    \vskip8pt
    \textbf{Today:} Mercury and Mars, the smaller siblings at opposite evolutionary limits.
  \end{columns}
\end{frame}


% ════════════════════════════════════════════════════════════
\sectionimage{mercury_caloris_basin}{Caloris Basin, Mercury. MESSENGER (NASA/JHUAPL)}
\section{Mercury Overview and Surface}
% ════════════════════════════════════════════════════════════

% ── Mercury at a glance ─────────────────────────────────────
\begin{frame}{Mercury at a Glance}
  \centering
  \small
  \begin{tabular}{l r l}
    \toprule
    \textbf{Parameter} & \textbf{Value} & \textbf{Note} \\
    \midrule
    Mass & $0.0553\,\Mearth$ & smallest planet \\
    Radius (mean) & $0.383\,\Rearth$ & $2440$ km \\
    Mean density & $5429$ kg\,m$^{-3}$ & uncompressed $\sim 5300$ kg\,m$^{-3}$ \\
    Semimajor axis & $0.387$ AU & innermost \\
    Orbital period & $87.97$ d & \\
    Rotation period & $58.65$ d & 3:2 spin-orbit resonance \\
    Eccentricity & $0.206$ & largest of any planet \\
    $T_{\mathrm{surf}}$ (day/night) & $700$/$100$ K & no atmosphere to redistribute \\
    Magnetic field & $\sim 1\%$ Earth's & weak but global \\
    \bottomrule
  \end{tabular}
  \vskip6pt
  \sourcenote{NASA Mercury Fact Sheet (NSSDC)}
\end{frame}

\begin{frame}
  \ipsherokey[0.66]{mercury_interior}{Mercury's internal structure from MESSENGER: crust, thin silicate mantle, liquid outer core and solid inner core. Credit: Course-original figure.}{A crust of 30--50 km and a 400 km mantle overlie a metallic core that fills 83\% of the radius.}
\end{frame}

% ── 3:2 resonance ─────────────────────────────────────────
\begin{frame}{The 3:2 Spin-Orbit Resonance}
  Mercury rotates exactly three times for every two orbits.
  \vskip8pt
  \begin{itemize}
    \item Discovered by radar (Pettengill \& Dyce 1965); stabilised by the high eccentricity ($e = 0.206$)
    \item The solar day lasts $176$ Earth days, exactly $2$ Mercury orbits
    \item The solar tidal torque on a permanent quadrupole moment locks the resonance
  \end{itemize}
\end{frame}

\begin{frame}
  \ipshero{messenger_mla_global}{MESSENGER MLA global topography: about 10 km of relief across ancient cratered highlands and younger volcanic plains (3.7--3.9 Ga). Credit: NASA/JHUAPL/CIW.}
\end{frame}

% ── Caloris Basin ─────────────────────────────────────────
\begin{frame}
  \ipsherokey[0.66]{mercury_caloris_basin}{Caloris Basin, about 1550 km across, formed about 3.8 Ga, with its concentric rings and lava-flooded floor. Credit: NASA/JHUAPL (MESSENGER).}{The largest well-preserved basin on Mercury and its chronostratigraphic reference; the focused seismic shock built the antipodal ``weird terrain''.}
\end{frame}

% ── Lobate scarps ─────────────────────────────────────────
\begin{frame}
  \ipsherokey[0.65]{mercury_lobate_scarp}{A lobate scarp: a compressional thrust fault with up to 3 km of relief and over 1000 km of length. Credit: NASA/JHUAPL (MESSENGER); Byrne et al.\ (2014).}{Secular cooling shrank Mercury by $\Delta R \approx 7$ km since the late bombardment; the area loss $\Delta A/A = 2\Delta R/R \approx 0.6\%$ is taken up by thrust scarps across one unbroken lid.}
\end{frame}

% ── Polar volatiles ───────────────────────────────────────
\begin{frame}
  \ipsherokey[0.65]{messenger_mla_polar}{North polar topography, illumination and temperature from MESSENGER MLA: the permanently shadowed craters coincide. Credit: Neumann et al.\ (2013).}{Crater floors below 100 K are cold traps; neutron spectrometry confirms water ice, delivered by comet and asteroid impacts, on the hottest planet.}
\end{frame}

% ── Surface chemistry ────────────────────────────────────
\begin{frame}
  \ipsherokey[0.66]{nittler2020_mercury_chemistry}{Surface composition of Mercury from the MESSENGER X-ray and gamma-ray spectrometers. Credit: Nittler et al.\ (2020).}{High S, K, Cl and Na despite 700 K days, FeO below 2\% and a high Mg/Si ratio: a highly reducing mantle formed under low oxygen fugacity.}
\end{frame}

% ── Exosphere and magnetosphere ──────────────────────────
\begin{frame}{Exosphere and Magnetosphere}
  A surface-bounded exosphere fed by sputtering, photon desorption and micrometeoroid vaporisation.
  \vskip8pt
  \begin{itemize}
    \item Surface pressure $\sim 10^{-12}$ bar; detected species include Na, K, Ca, Mg, He, H
    \item The sodium tail extends $\sim 10^6$ km anti-sunward under radiation pressure
    \item A miniature magnetosphere stands off the solar wind at $\sim 1.5\,R_{Hg}$ subsolar distance
  \end{itemize}
\end{frame}


% ════════════════════════════════════════════════════════════
\sectionimage{mercury_lobate_scarp}{Mercury surface detail. MESSENGER (NASA/JHUAPL)}
\section{Mercury Interior and Dynamo}
% ════════════════════════════════════════════════════════════

% ── Mercury interior layers ──────────────────────────────
\begin{frame}
  \ipsherokey[0.66]{margot2018_mercury_layers}{Mercury's layers from the libration and gravity measurements. Credit: Margot et al.\ (2018).}{A crust of $\sim 35$ km and a $\sim 400$ km silicate mantle above a liquid outer core below $\sim 2020$ km radius ($\sim 83\%$ of the planet) and a solid, possibly Fe-S, inner core: core mass fraction $\sim 74\%$.}
\end{frame}

% ── Margot libration ─────────────────────────────────────
\begin{frame}
  \ipsherokey[0.65]{margot2007_libration}{Radar speckle interferometry measures Mercury's 88-day forced libration. Credit: Margot et al.\ (2007), Fig.~3.}{The amplitude of $35.8 \pm 2$ arcseconds gives $C_m/C \approx 0.5$: only the silicate shell librates, so the mantle floats on a liquid core.}
\end{frame}

\begin{frame}
  \ipsherokey[0.6]{margot2007_libdata}{Monte Carlo distributions of the mantle's share of the moment of inertia, $C_m/C$, from the libration amplitude and $C_{22}$, for $C/MR^2$ of 0.325 (red) and 0.380 (blue). Reproduced from Margot et al.\ (2007).}{Every distribution peaks near $C_m/C \approx 0.5$, far from the value of 1 a fully solid Mercury requires; MESSENGER later gave $C/MR^2 = 0.346$, between the two cases.}
\end{frame}

% ── Why iron-rich ───────────────────────────────────────
\begin{frame}
  \ipsherokey{terrestrial_interiors_comparison}{Interiors of the terrestrial planets to scale: Mercury's core is out of proportion. Credit: NASA Solar System Exploration.}{\textbf{Giant impact} mantle stripping leads; \textbf{thermal vaporisation} is disfavoured by the retained volatiles (Na, S, K) and \textbf{selective condensation} of metal before silicate is no longer favoured.}
\end{frame}

% ── Franco mercury impact outcomes ──────────────────────
\begin{frame}
  \ipsherokey[0.6]{franco2022_mercury_outcomes}{Remnant mass against semimajor axis in N-body terrestrial-formation simulations for surface-density slopes $x = 0.5$--$5.5$; triangles mark the real planets. Credit: Franco et al.\ (2022), Fig.~A1.}{A Mercury-mass, iron-rich remnant at $0.39$ AU appears in $\ll 1\%$ of trials: a single-impact origin is dynamically rare.}
\end{frame}

% ── Wicht offset dipole ─────────────────────────────────
\begin{frame}
  \ipsherokey[0.61]{wicht_offset_dipole}{Mercury's dipole, offset $\sim 0.2\,R_{Hg}$ northward, in a dynamo model with a stably stratified layer. Credit: Wicht \& Heyner (2017).}{Convection in a thin outer-core shell gives a surface field of only $\sim 200$--$400$ nT (Earth $\sim 50,000$ nT), and the stratified liquid layer above suppresses the high-degree harmonics.}
\end{frame}


% ════════════════════════════════════════════════════════════
\sectionimage{mariner4_mars}{Mariner 4 view of Mars (1965). NASA/JPL}
\section{Mars Overview and Moons}
% ════════════════════════════════════════════════════════════

% ── Mars at a glance ──────────────────────────────────────
\begin{frame}{Mars at a Glance}
  \centering
  \small
  \begin{tabular}{l r l}
    \toprule
    \textbf{Parameter} & \textbf{Value} & \textbf{Note} \\
    \midrule
    Mass & $0.107\,\Mearth$ & one tenth of Earth \\
    Radius (mean) & $0.532\,\Rearth$ & $3389$ km \\
    Mean density & $3934$ kg\,m$^{-3}$ & lowest of the rocky planets \\
    Semimajor axis & $1.524$ AU & \\
    Orbital period & $686.97$ d & \\
    Rotation period & $24.62$ h & similar to Earth \\
    Axial tilt & $25.2^\circ$ & seasons like Earth's \\
    Eccentricity & $0.0934$ & strong climate forcing \\
    $T_{\mathrm{surf}}$ (day/seasonal) & $-130$ to $+30$\,$^\circ$C & \\
    Surface pressure & $\sim 6$ mbar & $0.6\%$ of Earth's \\
    Magnetic field & none global & local crustal anomalies \\
    \bottomrule
  \end{tabular}
  \vskip6pt
  \sourcenote{NASA Mars Fact Sheet (NSSDC)}
\end{frame}

\begin{frame}
  \ipshero{mariner4_mars}{The first close-up of Mars, Mariner 4 on 15 July 1965: cratered terrain ended the speculation about canals and vegetation. Credit: NASA/JPL.}
\end{frame}

% ── Phobos and Deimos ─────────────────────────────────────
\begin{frame}{Phobos and Deimos}
  Two small rubble-pile moons: captured carbonaceous asteroids, or re-accreted from a giant-impact debris disk.
  \vskip8pt
  \begin{itemize}
    \item \textbf{Phobos}: diameter $22$ km, orbital period $7.7$ h at $9376$ km radius
    \item \textbf{Deimos}: diameter $12$ km, orbital period $30.3$ h at $23463$ km radius
    \item Phobos orbits below the synchronous radius and will tidally disrupt in $30$--$50$ Myr
  \end{itemize}
\end{frame}

% ── Hyodo phobos impact ─────────────────────────────────
\begin{frame}
  \ipsherokey[0.61]{hyodo2017_phobos_impact}{SPH simulation of a Borealis-scale impact on Mars that forms an extended debris disk. Credit: Hyodo et al.\ (2017).}{Phobos and Deimos accrete from the disk, which reproduces their orbits and predicts a fraction of Mars-derived material in both moons; the composition (Mars mantle or D-type) is still debated.}
\end{frame}

% ── MMX timeline ─────────────────────────────────────────
\begin{frame}
  \ipsherokey[0.61]{mmx_timeline}{Operation plan of JAXA's Martian Moons eXploration: a three-year stay at Mars in five phases. Credit: Kuramoto et al.\ (2022), Fig.~4.}{Launch 2026, touchdowns on Phobos and a Deimos flyby between the Mars eclipses, sample return 2031: the samples will decide between capture and impact origin.}
\end{frame}


% ════════════════════════════════════════════════════════════
\sectionimage{terrestrial_planets}{Terrestrial planets to scale. Credit: NASA/JPL-Caltech}
\section{Mars' Interior: The InSight Revolution}
% ════════════════════════════════════════════════════════════

% ── Stahler interior structure ──────────────────────────
\begin{frame}
  \ipsherokey[0.56]{stahler2021_mars_structure}{Mars' interior from InSight SEIS: crust, mantle and core. Credit: St\"ahler / Khan / Knapmeyer-Endrun (2021).}{Crust $24$--$72$ km, mantle to $\sim 1560$ km with the olivine-wadsleyite transition near $1000$--$1100$ km, a liquid core boundary read at $\sim 1830$ km in 2021; the 2023 reanalyses place a molten silicate layer above a smaller core ($\sim 1650$--$1675$ km, $\sim 6500$ kg\,m$^{-3}$).}
\end{frame}

% ── Marsquakes ───────────────────────────────────────────
\begin{frame}
  \ipsherokey{stahler2021_marsquakes}{Marsquake waveforms recorded by InSight: P and S arrivals and core-reflected ScS phases. Credit: St\"ahler et al.\ (2021).}{ScS shear reflections confirm a liquid core; far-side quakes and impacts (2023) refine its radius to $\sim 1650$--$1675$ km; over $1300$ events were recorded, the largest S1222a at $M = 4.7$.}
\end{frame}

% ── Plesa convection ─────────────────────────────────────
\begin{frame}
  \ipsherokey{plesa2022_convection}{A 3D thermal-evolution model of Mars' mantle consistent with the InSight constraints. Credit: Ple\v{s}a et al.\ (2022).}{A few long-wavelength upwellings and sluggish lower-mantle flow under a thick stagnant lid; the Tharsis plume persists to the present.}
\end{frame}

\begin{frame}
  \ipshero{plesa2022_crustalthickness}{Present-day crustal thickness of Mars: the global mean spans 40.6 km (31 km at InSight) in the thin end-member to 71.4 km (49 km at InSight) in the thick one. Ple\v{s}a et al.\ (2022).}
\end{frame}


% ════════════════════════════════════════════════════════════
\sectionimage{olympus_mons}{Olympus Mons. NASA/JPL/Mars Express (ESA)}
\section{Mars Geology and Surface Highlights}
% ════════════════════════════════════════════════════════════

% ── Tanaka geological periods ────────────────────────────
\begin{frame}
  \ipsherokey{tanaka2014_periods}{The three classical periods of Mars' geological time scale. Credit: Tanaka et al.\ (2014).}{\textbf{Noachian} ($> 3.7$ Ga): heavy bombardment, valley networks, warm and wet episodes; \textbf{Hesperian} ($3.7$--$3.0$ Ga): outflow channels and large-scale volcanism; \textbf{Amazonian} ($< 3.0$ Ga): cold, hyper-arid, thin atmosphere.}
\end{frame}

% ── Tanaka geological map ───────────────────────────────
\begin{frame}
  \ipsherokey{tanaka2014_geomap}{The global geological map of Mars. Credit: Tanaka et al.\ (2014).}{Tharsis dominates the western hemisphere, Hellas, Argyre and Isidis are the giant impact basins, and the northern lowlands are young, smooth and partly buried.}
\end{frame}

\begin{frame}
  \ipshero{olympus_mons}{Olympus Mons: a shield volcano $\sim 21$ km high and $600$ km wide, built by stationary hotspot volcanism over a stagnant lid. Credit: ESA/DLR/FU Berlin.}
\end{frame}

\begin{frame}
  \ipshero{valles_marineris}{Valles Marineris: a graben system $4000$ km long and up to $7$ km deep, opened by crustal extension under the Tharsis load. Credit: NASA/JPL-Caltech.}
\end{frame}

% ── Valley networks (water evidence) ────────────────────
\begin{frame}
  \ipsherokey[0.66]{mars_valley_networks_viking}{Dendritic valley networks carved into Noachian terrain, imaged by the Viking Orbiters. Credit: NASA/JPL/USGS.}{Confluences, meanders and tributaries confirm liquid-water erosion, with drainage densities near terrestrial arid networks, active for $> 100$ Myr in the late Noachian.}
\end{frame}

% ── Kite valley distribution ────────────────────────────
\begin{frame}
  \ipsherokey[0.66]{kite2022_valley_distribution}{Late Noachian valley networks ($\sim 3.6$ Ga and older, top) against late Hesperian to Amazonian fans ($3.5$--$3$ Ga, bottom); elevations $-6$ to $+6$ km. Credit: Kite et al.\ (2022), Fig.~1.}{The early valleys sit on high ground and the later fans in a mid-latitude band: fluvial control shifts from \textbf{elevation} to \textbf{latitude} as the climate cools.}
\end{frame}

% ── Bibring mineralogy timeline ─────────────────────────
\begin{frame}
  \ipsherokey[0.61]{bibring2006_timeline}{The OMEGA mineralogy time line: Phyllosian, Theiikian, Siderikian. Credit: Bibring et al.\ (2006), Fig.~5.}{\textbf{Phyllosian}: phyllosilicates from neutral-pH water; \textbf{Theiikian}: sulfates from acidic, evaporating water; \textbf{Siderikian} (from $\sim 3.5$ Ga): anhydrous ferric oxides under hyper-arid conditions.}
\end{frame}

% ── Bibring global map ───────────────────────────────────
\begin{frame}
  \ipsherokey[0.65]{bibring2006_globalmap}{Global distribution of phyllosilicates (red), sulfates (blue) and other hydrated minerals (yellow). Credit: Bibring et al.\ (2006), Fig.~3.}{Clays cluster in the ancient southern Noachian highlands, sulfates in the equatorial evaporative basins.}
\end{frame}

% ── Ehlmann olivine ────────────────────────────────────
\begin{frame}
  \ipsherokey{ehlmann2014_olivine}{Widespread unweathered olivine in the Martian crust. Credit: Ehlmann \& Edwards (2014).}{Olivine alters quickly in liquid water to serpentine and smectite clays, so its survival rules out a persistent global ocean: water activity was episodic or localised.}
\end{frame}

\begin{frame}
  \ipshero[0.76]{ehlmann2014_spectra}{Mars surface compositions from infrared remote sensing and rover instruments: basalts, hydrated silicates, sulfates, carbonates, hematite and opaline silica record a wide range of aqueous and igneous environments. Ehlmann \& Edwards (2014).}
\end{frame}

\begin{frame}
  \ipshero[0.76]{msr_architecture}{Mars Sample Return concept: cache at Jezero, retrieval lander and ascent vehicle, orbital capture, Earth entry and laboratory analysis; under rebaselining after the 2024 cost review. Course figure.}
\end{frame}


% ── Break ───────────────────────────────────────────────────
\breakslide


% ════════════════════════════════════════════════════════════
\sectionimage{family_portrait}{Solar system, Voyager 1 family portrait. Credit: NASA/JPL-Caltech}
\section{Blackboard Derivation: Jeans Escape Flux}
% ════════════════════════════════════════════════════════════

% ── Setup: goal and strategy ──────────────────────────────
\begin{frame}{Goal and strategy}
  \begin{columns}[c]
    \column{\recapfigcolnarrow}
    \ipsrecapfig[0.62\textheight]{board_sketch_l10}
    \column{\recaptextcolwide}
    \small
    \textbf{Goal:} the rate at which thermal energy alone lets molecules escape, applied to Mars.
    \vskip4pt
    \begin{itemize}\setlength\itemsep{2pt}
      \item \textbf{The exobase:} mean free path equals scale height; above it, ballistic flight
      \item $v_{\mathrm{th}} = \sqrt{2 \kB T / m}$ against $v_{\mathrm{esc}} = \sqrt{2 G M / r_{\mathrm{exo}}}$
      \item Count the molecules in the fast tail that point outward
    \end{itemize}
    \vskip4pt
    \keyresult{A molecule faster than $v_{\mathrm{esc}}$ leaves; a slower one falls back.}
  \end{columns}
\end{frame}

\begin{frame}{The Escape Parameter}
  \keyresult{$$\lambda \equiv \left(\frac{v_{\mathrm{esc}}}{v_{\mathrm{th}}}\right)^{\!2} = \frac{G\,M\,m}{\kB\,T\,r_{\mathrm{exo}}}$$}
  \vskip6pt
  \begin{itemize}
    \item $\lambda \gg 1$: molecules gravitationally bound; escape rare
    \item $\lambda \lesssim 1$: thermal energy comparable to escape energy; rapid loss
    \item Depends strongly on molecular mass $m$; independent of magnetic field, ionisation and solar wind
  \end{itemize}
\end{frame}

% ── Jeans flux derivation ────────────────────────────────
\begin{frame}{Result: The Jeans Formula}
  Integrating upward velocity over the upper hemisphere with cutoff $v \geq v_{\mathrm{esc}}$:
  $$\Phi_J = n\,\langle v\rangle\,\frac{1}{4}\,(1+\lambda)\,e^{-\lambda}$$
  \vskip4pt
  where $\langle v\rangle = \sqrt{8 \kB T / (\pi m)}$ is the mean thermal speed.
  \vskip6pt
  Substituting and simplifying:
  \keyresult{$$\Phi_J \;=\; n\,\sqrt{\frac{\kB T}{2\pi m}}\,(1+\lambda)\,e^{-\lambda}$$}
  \vskip4pt
  Jeans (1925). Limits: $\lambda \to 0$ gives the effusion flux $n\sqrt{\kB T/(2\pi m)}$; for $\lambda \gg 1$ the flux falls as $\lambda e^{-\lambda}$, so $\lambda$ from $5$ to $10$ cuts it by $\sim 80$.
\end{frame}

\begin{frame}
  \ipshero[0.76]{maxwell_boltzmann_jeans}{Maxwell-Boltzmann speed distributions of atomic hydrogen at the Earth (1000 K) and Mars (270 K) exobases: only the tail above the escape velocity leaves. Credit: Course-original figure.}
\end{frame}

% ── Mars application ─────────────────────────────────────
\begin{frame}{Apply to Mars}
  $M = 6.4\times 10^{23}$ kg, $r_{\mathrm{exo}} \approx 3.6\times 10^6$ m, $T_{\mathrm{exo}} \approx 270$ K.
  \vskip2pt
  Escape velocity at the exobase: $v_{\mathrm{esc}} \approx 4.9\,\mathrm{km/s}$
  \vskip4pt
  \centering
  \small
  \begin{tabular}{l c c}
    \toprule
    \textbf{Species} & $\lambda$ & $e^{-\lambda}$ \\
    \midrule
    H & $5.3$ & $5\times 10^{-3}$ \\
    $\mathrm{H_2}$ & $10.6$ & $2$--$3\times 10^{-5}$ \\
    He & $21.2$ & $\sim 10^{-9}$ \\
    O & $84.8$ & $10^{-37}$ \\
    $\mathrm{CO_2}$ & $230$ & $10^{-100}$ \\
    \bottomrule
  \end{tabular}
  \vskip4pt
  \keyresult{Mars cannot lose $\mathrm{CO_2}$ thermally. Only H and $\mathrm{H_2}$ escape via Jeans; heavier species need non-thermal channels.}
\end{frame}

\begin{frame}
  \ipshero[0.76]{mars_dust_storm}{Mars before and during the 2018 global dust storm, imaged by the Mars Reconnaissance Orbiter: lofted dust hid the surface and blocked solar power. Credit: NASA/JPL-Caltech/MSSS, public domain.}
\end{frame}


% ════════════════════════════════════════════════════════════
\sectionimage{mars_valley_networks_viking}{Dendritic channel networks in the cratered Martian highlands, imaged by the Viking Orbiters. Credit: NASA/JPL/USGS}
\section{Mars Atmosphere and Climate Loss}
% ════════════════════════════════════════════════════════════

% ── Early Mars climate puzzle & cold baseline ───────────
\begin{frame}
  \ipsherokey[0.6]{early_mars_climate_schematic}{The early Mars climate problem: a faint young Sun against valley networks that need liquid water. Credit: Wordsworth (2016, 2017, 2021); Kite \& Conway (2024).}{At $4$ Ga the Sun gave $\sim 75\%$ of its modern flux and a pure $\mathrm{CO_2}$ atmosphere yields $T \le 225$ K; transient $\mathrm{H_2}$-$\mathrm{CO_2}$ collision-induced absorption triggers snowmelt excursions of $10^5$ to $10^6$ yr.}
\end{frame}

% ── Wordsworth phase diagram ────────────────────────────
\begin{frame}
  \ipsherokey[0.6]{wordsworth2016_phasediagram}{Long-term climate regimes of Mars: mean surface temperature (divider $\sim 280$ K) against surface $\mathrm{H_2O}$ inventory (divider $\sim 200$ m GEL), with cross-sections. Credit: Wordsworth (2016), Fig.~7.}{The four quadrants separate ice-capped highlands from lowland water; the geomorphology favours a cold, dry baseline with episodic melting.}
\end{frame}

% ── Hu carbon evolution ─────────────────────────────────
\begin{frame}
  \ipsherokey[0.6]{hu2015_carbon_evolution}{A coupled atmosphere-surface-interior model of Mars' carbon inventory through time; initial $\mathrm{CO_2}$ of $0.32$--$0.50$ bar (red) and $0.25$--$0.29$ bar (blue). Credit: Hu et al.\ (2015).}{Late Noachian carbon is reconstructed at $\sim 0.5$--$1.8$ bar, then lost to atmospheric escape and carbonate burial: the modern $6$ mbar is the remnant of the primordial inventory.}
\end{frame}

% ── Kite schematic ──────────────────────────────────────
\begin{frame}
  \ipsherokey[0.65]{kite2022_schematic}{Climate state against $\mathrm{CO_2}$ loss and lost non-$\mathrm{CO_2}$ greenhouse forcing, from the high early valleys to the later low fans. Credit: Kite et al.\ (2022), Fig.~6.}{The transition cooled Mars by $\gtrsim 10$ K, driven primarily by \textbf{waning non-$\mathrm{CO_2}$ forcing} rather than by the thinning of the $\mathrm{CO_2}$ atmosphere.}
\end{frame}

% ── Jakosky loss ─────────────────────────────────────────
\begin{frame}
  \ipsherokey[0.62]{jakosky2018_loss}{MAVEN oxygen escape channels extrapolated back to $\sim 3.5$ Ga. Credit: Jakosky et al.\ (2018).}{\textbf{Dissociative recombination} dominates modern O loss, \textbf{pick-up ions} show the steepest historical decline, \textbf{ion outflow} and solar-wind \textbf{sputtering} add the rest; the present total escape is $\sim 2$ kg\,s$^{-1}$, and the heavy species require these non-thermal channels.}
\end{frame}

\begin{frame}
  \ipshero[0.76]{jakosky2018_hloss}{Hydrogen corona column density at Mars over a Mars year: the escape varies by an order of magnitude and peaks near perihelion, when water vapour reaches the altitudes where it is photolysed. Jakosky et al.\ (2018).}
\end{frame}


% ════════════════════════════════════════════════════════════
\sectionimage{mariner4_mars}{Early Mars from Mariner 4 (1965). Credit: NASA/JPL}
\section{Mars Magnetism}
% ════════════════════════════════════════════════════════════

% ── Acuna magnetic map ──────────────────────────────────
\begin{frame}
  \ipsherokey[0.66]{acuna1999_magmap}{Crustal magnetic field of Mars from Mars Global Surveyor. Credit: Acu\~na et al.\ (1999).}{No global dipole today, but intense remnant magnetisation in the southern highlands requires a vigorous early dynamo that shut down; the northern lowlands and the major basins were demagnetised by impacts.}
\end{frame}

% ── Acuna dipoles ───────────────────────────────────────
\begin{frame}
  \ipsherokey[0.61]{acuna1999_dipoles}{Polar view of $B_r$; the companion equatorial map shows east-west lineated anomalies at mid-latitudes. Credit: Acu\~na et al.\ (1999).}{Basaltic crust recorded the field as it cooled, possibly through reversals; the pre-Borealis crust preserves the primordial dynamo, whose cessation is dated to $\sim 4.1$ Ga (refined to $\sim 3.7$ Ga).}
\end{frame}

\begin{frame}
  \ipsherokey[0.56]{dynamo_lifetimes}{Dynamo lifetimes of the rocky planets: Earth active, Mercury weak in a thin liquid shell, Mars for its first 500 to 800 Myr with crustal fields as young as about 3.7 Ga, Venus with no detectable field. Course figure.}{A magnetic shield suppresses ion escape by up to an order of magnitude over Gyr; an unmagnetised Mars-sized planet still loses its atmosphere on $10^9$-year timescales.}
\end{frame}


% ════════════════════════════════════════════════════════════
\sectionimage{terrestrial_planets}{Terrestrial planets to scale. Credit: NASA/JPL-Caltech}
\section{Comparative Payoff: Mercury vs Mars}
% ════════════════════════════════════════════════════════════

% ── BepiColombo hero ────────────────────────────────────────
\begin{frame}
  \ipshero[0.76]{wicht_magnetosphere}{Mercury's magnetosphere, the smallest in the solar system, with its subsolar magnetopause at $\sim 1.5\,R_{Hg}$; BepiColombo (ESA/JAXA) enters orbit in late 2026 with two orbiters to map the offset field, the core and the exosphere. Credit: Wicht \& Heyner (2017).}
\end{frame}

% ── Size and distance ─────────────────────────────────────
\begin{frame}{Size and Distance Set the Trajectory}
  \begin{itemize}\setlength\itemsep{4pt}
    \item \textbf{Small bodies cool fast}: surface-area-to-volume ratio $\propto 1/R$; Mars cooled faster than Earth, so its dynamo shut off early
    \item \textbf{Mercury keeps a weak dynamo}: light elements (S, Si) keep a thin outer-core shell liquid; thermal stratification gives the weak, offset field
    \item \textbf{Lesson}: dynamo persistence needs (a) core size, (b) a slow-cooling mantle, (c) chemical buoyancy from inner-core growth; Mars failed (b), Mercury keeps (a), Earth has all three
  \end{itemize}
\end{frame}

\begin{frame}
  \ipsherokey[0.65]{size_distance_regimes}{Size and distance set the trajectory (schematic); the shaded regime boundary is qualitative. Course figure.}{Larger bodies keep their atmospheres, and their dynamos, for longer; size and distance govern core cooling, dynamos and volatile retention.}
\end{frame}

\begin{frame}{Summary}
  \begin{enumerate}\setlength\itemsep{6pt}
    \item \textbf{Mercury}: 3:2 spin-orbit resonance, a $\sim 74\%$ iron core from mantle stripping, lobate scarps from $\Delta R \approx 7$ km of \textbf{contraction}, a weak offset dynamo
    \item \textbf{Mars}: a liquid core of $\sim 1650$--$1675$ km below a molten silicate layer (2023 reanalyses), a wet Noachian record, and a dynamo that died between $\sim 4.1$ and $\sim 3.7$ Ga
    \item \textbf{Jeans escape} $\Phi_J \propto (1+\lambda)e^{-\lambda}$ lets H go but keeps $\mathrm{CO_2}$; \textbf{non-thermal loss} (sputtering, ion pickup) removes the heavy species
  \end{enumerate}
  \vskip6pt
  \textbf{Next: Lecture 11}
\end{frame}

% ── Final slide ─────────────────────────────────────────────
\begin{frame}[plain,noframenumbering]
  \begin{tikzpicture}[remember picture, overlay]
    \ipscontain{title_background}
    \fill[ipsPrimary, opacity=0.70] (current page.north west) rectangle (current page.south east);
  \end{tikzpicture}
  \vfill
  \begin{center}
    {\color{white}\Large\bfseries Questions?}
    \vskip20pt
    {\color{ipsLightFill}\normalsize Next: \textbf{Lecture~11. Gas \& Ice Giants}}
    \vskip16pt
    {\color{ipsLightFill!70}\small Lecture notes: \texttt{formingworlds.github.io/IntroductionToPlanetaryScience}}
  \end{center}
  \vfill
\end{frame}

\end{document}
